\documentclass[11pt,a4paper]{article}
\usepackage[T1]{fontenc}\usepackage[margin=26mm]{geometry}
\usepackage{amsmath,amssymb,amsthm,mathtools,booktabs,lmodern,microtype}
\usepackage[hidelinks]{hyperref}\usepackage{xurl}
\hypersetup{pdftitle={Positivity at the First Feedback Boundary in Thue Morse Pressure},pdfauthor={Research note prepared for Vladimir Reshetnikov with OpenAI}}
\setlength{\parskip}{3pt}
\newtheorem{theorem}{Theorem}[section]\newtheorem{lemma}[theorem]{Lemma}
\newtheorem{proposition}[theorem]{Proposition}\newtheorem{corollary}[theorem]{Corollary}
% ed. (2026-10-01): unnumbered environment for ProveIt editorial notes (no counter changes).
\theoremstyle{remark}\newtheorem*{ednote}{Editorial note (ProveIt, 2026-10-01)}
\newcommand{\Z}{\mathbb Z}\newcommand{\ii}{\mathrm i}\newcommand{\ev}{\operatorname{ev}}
\title{Positivity at the First Feedback Boundary in Thue Morse Pressure}
\author{Research note prepared for Vladimir Reshetnikov with OpenAI}\date{1 October 2026}
\begin{document}\maketitle
\begin{abstract}
We study the coefficient of degree $4m$ in integer Thue--Morse pressure, where nonlinear eigenvalue feedback first enters. An exact regularization separates the boundary eigenfunction response into a strictly positive finite resolvent and an absolutely convergent mixed inverse-orbit sum. A strict dyadic tangent comparison and a lattice local limit argument give a positive diagonal asymptotic. An explicit analytic estimate proves positivity for every $m\ge70$, and exact integer enclosure certificates cover $2\le m\le69$. Together they prove that the degree-$4m$ pressure coefficient is positive for every integer $m\ge2$. This continues the positive triangle strictly below degree $4m$ without changing the earlier reports.
\end{abstract}
\paragraph{Status.}
This is unrefereed ordinary mathematics with exact arithmetic certificates, not a Lean formalization. No worldwide priority claim or general formula for the first later negative coefficient is made. The proofs distinguish the infinite-order estimates from their finite verification complement.

\section{The boundary result}
Use the doubling-map pressure $p_m(c)$ of $m\log\cos^2\!\pi(x-c)$ in the sum-over-preimages convention. Set $a=\tan t$ and
\[
 (V_af)(x)=\sum_{2y=x\bmod1}(\cos\pi y+a\sin\pi y)^{2m}f(y).
\]
Let $h_a(0)=1$ normalize its local analytic Perron eigenfunction. Write
\[
 R_m=\frac{2(2^{2m}-1)\zeta(2m)}{\pi^{2m}},\qquad
 H_{m,r}=[a^{2r}]h_a(1/2),\qquad E_{m,r}=[t^{2m+2r}]p_m(t/\pi).
\]
The preceding companion proves $H_{m,r},E_{m,r}>0$ whenever $1\le r<m$ \cite{triangle}. Here the new response is $H_{m,m}$ and the new pressure coefficient is $E_{m,m}=[t^{4m}]p_m(t/\pi)$.

\begin{theorem}\label{thm:main}
For every integer $m\ge2$,
\[
 H_{m,m}>0,\qquad [t^{4m}]p_m(t/\pi)>0.
\]
Moreover, put $b_j=\tan(\pi/2^{j+1})$ and $B(a)=\prod_{j\ge1}(1+ab_j)$. There is a unique $a_*>0$ satisfying
\[
 \sum_{j\ge1}\frac{a_*b_j}{1+a_*b_j}=1.
\]
With
\[
 C=\frac{B(a_*)}{a_*},\qquad p_j=\frac{a_*b_j}{1+a_*b_j},\qquad
 \sigma^2=\sum_{j\ge1}p_j(1-p_j),
\]
we have the asymptotic
\begin{equation}\label{eq:asymptotic}
 H_{m,m}\sim\frac{2(2C/\pi)^{2m}}{\sqrt{4\pi m\sigma^2}}
 \qquad(m\to\infty).
\end{equation}
\end{theorem}
The constants are defined by convergent products and sums; approximate values are $a_*\approx0.822918569054$, $C\approx4.046860022$, and $\sigma^2\approx0.7041449521$. These decimals are orientation only and are not used in the sign proof. The finite-complement certification is described in Section~\ref{sec:finite}. Combined with the earlier triangle, this proves $H_{m,r},E_{m,r}>0$ for every $1\le r\le m$. The response range is sharp as a universal statement: $H_{2,3}=-488/27<0$.
% ed. (2026-10-01): the pressure statement is re-proved in the full range; the response statements stay here
\begin{ednote}
The pressure statement $E_{m,m}>0$ is re-proved, by a separate large-order bound and separate
finite certificates, as the case $r=m$ of Theorem~1.1 of \path{../Thue_Morse_Integer_Pressure_Full_Range/}
($E_{m,r}>0$ for every $m\ge2$ and $1\le r<2m$), so it does not depend on the trace archives
discussed in the note in Section~\ref{sec:finite}. The response statement $H_{m,m}>0$ and the
asymptotic \eqref{eq:asymptotic} are this note's own; \path{../Thue_Morse_Integer_Pressure_Higher_Positive/} had
proved the case $m=2$ ($H_{2,2}=536/27$), and \path{../Thue_Morse_Integer_Pressure_Beyond_Boundary/} extends the
asymptotic to every fixed offset $s$ (the case $s=0$ is \eqref{eq:asymptotic}).
\end{ednote}

\section{The exact boundary regularization}
Put $d=2m$, $A=V_0$, and $h=F_{0,m}$. The atomic basis is
\[
 F_{j,m}(x)=\sin^{2m}(\pi x)\sum_{n\in\Z}(\pi(x+n))^{-(2m-j)},
 \quad0\le j\le2m-2.
\]
It spans the invariant Fourier space with indices $1-m,\ldots,m-1$ and satisfies $AF_{j,m}=2^{-j}F_{j,m}$. In particular $h(0)=1$, $h(1/2)=R_m$, and the rank-one spectral projection is $Pf=f(0)h$. These facts follow by splitting the periodization into even and odd integers; the Fourier model and its local pressure interpretation are reviewed in \cite{integer,triangle}. Let $Q=I-P$ and let $\ev f=f(1/2)$.

The eigenvalue and pressure satisfy
\begin{equation}\label{eq:feedback}
 \lambda_m(a)=1+a^dh_a(1/2),\qquad
 p_m(t/\pi)=2m\log\cos t+\log\lambda_m(\tan t).
\end{equation}
At degree $d$, the raw inverse-orbit coefficient no longer has the absolute-convergence estimate used below the boundary. Its exact regularization is as follows.

\begin{lemma}\label{lem:regularization}
If $C_N=[a^d](V_a^Nh)(1/2)$, then
\begin{equation}\label{eq:renormalized}
 H_{m,m}=\lim_{N\to\infty}(C_N-NR_m^2).
\end{equation}
For every finite $N$,
\begin{equation}\label{eq:finiteorbit}
 (V_a^Nh)(1/2)=\sum_{w\in\Z+1/2}(\pi w)^{-d}
       \prod_{j=1}^N\left(1+a\tan\frac{\pi w}{2^j}\right)^d.
\end{equation}
\end{lemma}
\begin{proof}
Evaluation at zero gives $(V_af)(0)=f(0)+a^df(1/2)$. Since $A^Nh=h$,
\[
 (V_a^Nh)(0)=1+NR_ma^d+O(a^{d+1}).
\]
The normalized iterates converge locally uniformly analytically to $h_a$, by the simple dominant eigenvalue and finite-dimensional spectral projection. Dividing by the last display proves~\eqref{eq:renormalized} by Cauchy's coefficient formula.

For~\eqref{eq:finiteorbit}, first use centered inverse representatives $|w|<2^{N-1}$. The cosine product contributes $(2^N\sin(\pi w/2^N))^{-d}$. Substituting the periodization defining $h(w/2^N)$ cancels its sine factor and unfolds the representatives into all half-integers. The finite tangent product is periodic in $w$ with period $2^N$ and has no pole there, so the resulting series is absolutely convergent.
\end{proof}

Let $t_j(w)=\tan(\pi w/2^j)$ and define
\[
 K_m(w)=[a^d]\prod_{j\ge1}(1+at_j(w))^d-\sum_{j\ge1}t_j(w)^d.
\]
The subtracted monomials choose all $d$ powers from just one tangent factor. Every remaining monomial uses at least two factors, with each selected exponent below $d$.

\begin{proposition}\label{prop:mixed}
The mixed orbit series is absolutely convergent, and
\begin{equation}\label{eq:decomposition}
 H_{m,m}=P_m+\sum_{w\in\Z+1/2}(\pi w)^{-d}K_m(w),
\end{equation}
where the finite resolvent is
\[
 P_m=\sum_{n\ge0}\bigl(\ev A^nV_dh-R_m^2\bigr)
     =\ev(I-A)^{-1}QV_dh.
\]
Here $V_a=\sum_{k=0}^da^kV_k$, and the inverse is restricted to $Q$.
\end{proposition}
\begin{proof}
Let $g(x)=\sin\pi x$ on $[0,1]$, extended periodically, and let $W_k$ replace each weight of $V_k$ by its absolute value. With $s=\sin(\pi x/2)$ and $c=\cos(\pi x/2)$,
\[
 Ag=sc(c^{d-1}+s^{d-1})\le\tfrac12g\qquad(d\ge4).
\]
For $1\le k<d$ and bounded $f$,
\begin{equation}\label{eq:vanishing}
 W_k|f|\le\tfrac12\binom dk\|f\|_\infty g.
\end{equation}
Indeed, after factoring $sc$ from $c^{d-k}s^k+s^{d-k}c^k$, the remaining sum is at most $c^{d-2}+s^{d-2}\le1$, by weighted AM--GM.

Fix a composition $k_1+\cdots+k_s=d$ with $s\ge2$ and all $k_i\ge1$. Its absolute insertion words have the form
\[
 \ev A^{n_0}W_{k_1}A^{n_1}\cdots A^{n_{s-1}}W_{k_s}h.
\]
The final power of $A$ after the last insertion fixes $h$. Inequality~\eqref{eq:vanishing} and $Ag\le g/2$ bound each word by a constant times $2^{-\sum n_i}$. Summing all zero-degree gaps is finite, and there are only finitely many compositions of $d$. Tonelli's theorem therefore proves absolute convergence of the mixed orbit monomials, including convergence of their finite truncations.

The pure terms in $C_N$ are exactly $\sum_{n=0}^{N-1}\ev A^nV_dh$. As $(V_dh)(0)=R_m$, their spectral projection is $R_mh$, giving the secular part $NR_m^2$. Subtraction in~\eqref{eq:renormalized} yields~\eqref{eq:decomposition}.
\end{proof}

\section{A strictly positive pure correction}
\begin{proposition}\label{prop:pure}
For every $m\ge2$,
\begin{equation}\label{eq:purebound}
 R_m(1-R_m)<P_m\le\frac43R_m(1-R_m).
\end{equation}
The upper equality holds only at $m=2$.
\end{proposition}
\begin{proof}
Write $F_{0,m}(x+1/2)=q_m(u)$, where $u=\sin^2\pi x$. The polynomial $q_m$ has degree $m-1$ and all coefficients strictly positive. A direct induction proves this: $q_1=1$, and twice differentiating the periodized secant expression gives
\begin{align}\label{eq:qrecurrence}
 q_{m+1}(u)=\frac1{2m(2m+1)}\bigl[&4u(1-u)^2q_m''(u)\notag\\
 &+2(1-u)(1+(4m-2)u)q_m'(u)
 +2m(1+2mu)q_m(u)\bigr].
\end{align}
Applied to $u^k$, its numerator has coefficients
\[
 2k(2k-1),\qquad 2m+2k(4m+1-4k),\qquad4(m-k)^2
\]
at $u^{k-1},u^k,u^{k+1}$. They are nonnegative for $0\le k\le m-1$ and propagate strict positivity. Formula~\eqref{eq:qrecurrence} follows from differentiating $(1-u)^{-m}q_m(u)$ and using $u'=2\sin z\cos z$, $u''=2(1-2u)$, with $z=\pi x$. Also $q_m(0)=R_m$ and $q_m(1)=1$, proving $R_m\le h\le1$.

The monomials have positive atomic expansions
\[
 u^k=\sum_{j=k}^{m-1}\beta_{m-k,j-k}F_{2j,m}(x),\qquad
 \beta_{p,\ell}=[z^{2\ell}](z/\sin z)^{2p}>0.
\]
This is the principal-part decomposition of $\csc^{2(m-k)}z$ after division by $\sin^{2m}z$. The remaining entire periodic difference vanishes at imaginary infinity. Positivity of $\beta_{p,\ell}$ follows from the Euler product for $z/\sin z$.
Thus
\[
 F_{0,m}(x+1/2)=\sum_{j=0}^{m-1}d_{m,j}F_{2j,m}(x),
 \qquad d_{m,j}>0,\quad d_{m,0}=R_m.
\]
The exact identity $V_d=AT_{1/2}$ now gives
\[
 P_m=\sum_{j=1}^{m-1}\frac{d_{m,j}R_{m-j}}{4^j-1}>0.
\]
At $x=1/2$ the sine and cosine weights agree, so $\ev V_dh=\ev Ah=R_m$. Consequently
\[
 \sum_{j=1}^{m-1}4^{-j}d_{m,j}R_{m-j}=R_m-R_m^2.
\]
Comparing termwise with $1<(1-4^{-j})^{-1}\le4/3$ proves~\eqref{eq:purebound} and its equality statement.
\end{proof}

\section{A strict tangent comparison and an exponential tail}
For positive odd $n$, put $x_j(n)=|\tan(n\pi/2^{j+1})|/n$. The tangent theorem in \cite{triangle} proves that its decreasing rearrangement is weakly log-majorized by $(b_1,b_2,\ldots)$, where $b_1=1$. We need the following strict refinement.

\begin{lemma}\label{lem:strict}
For every odd $n\ge3$, $x(n)$ is weakly log-majorized by
\[
 (\theta,b_2,b_3,\ldots),\qquad\theta=\frac{1+\sqrt2}{3}<1.
\]
Hence, for every $a>0$,
\begin{equation}\label{eq:strictproduct}
 B_{x(n)}(a):=\prod_j(1+ax_j(n))
 \le\frac{1+\theta a}{1+a}B(a).
\end{equation}
\end{lemma}
\begin{proof}
We recall the finite structure of the earlier tangent proof. Set $L=\lfloor\log_2n\rfloor$ and $s_j=|\sin(n\pi/2^{j+1})|/(n\sin(\pi/2^{j+1}))$. The sorted ratios $x_j^\downarrow/b_j$ cross $1$ only once, from below to above. After that crossing, the largest $k$ terms are exactly the first $k$ original terms, with product ratio
\[
 s_k^2\prod_{j=2}^{k-1}s_j\quad(k\ge2).
\]
These statements follow from the initial cotangent bounds and the tail counting argument in \cite{triangle}; all $s_j$ lie in $(0,1]$.

The largest normalized tangent is at most $\theta$. For $n=3$, its value is $\tan(3\pi/8)/3=\theta$. For $n\ge5$, write $q=2^{L+1}\ge8$. The original early terms are below $2/\pi$ and the late terms below $1/n$; the only remaining term obeys
\[
 x_{L+1}\le\frac{\cot(\pi/(2q))}{q-1}
 <\frac{2q}{\pi(q-1)}\le\frac{16}{7\pi}<\theta.
\]
Also $s_2\le\theta$: it is equal to $\theta$ for $n=3$, and for $n\ge5$ it is below $1/(5\sin(\pi/8))<4/5<\theta$.

Before the sorted crossing, the first ratio is at most $\theta$ and every later ratio is at most $1$. After the crossing, the original-prefix identity is at most $\theta$ because it contains $s_2$, with the case $k=2$ giving $s_2^2$. Every sorted prefix product is therefore at most $\theta$ times its $b$-prefix product. The comparison sequence is decreasing because $\theta>b_2$.

Weak log-majorization implies all elementary-symmetric coefficient inequalities: apply the increasing symmetric convex function $z\mapsto\sum_{|I|=r}\exp(\sum_{i\in I}z_i)$ to the logarithms of finite prefixes. This elementary majorization step is proved in \cite{triangle}. Geometric tails justify the monotone limit to infinite products and give~\eqref{eq:strictproduct}.
\end{proof}

Define the mixed central coefficient
\[
 M_d(x)=[a^d]\prod_j(1+ax_j)^d-\sum_jx_j^d.
\]
For a nonnegative summable sequence $x$ and any $a>0$, marking two distinct factor indices gives
\begin{equation}\label{eq:marked}
 M_d(x)\le\frac{d^2}{d-1}a^{2-d}e_2(x)B_x(a)^d.
\end{equation}
Indeed, a mixed degree-$d$ monomial with exponents $k_j$ is marked $\sum_{i<j}k_ik_j\ge d-1$ times. The marked generating function is
\[
 d^2a^2 B_x(a)^d\sum_{i<j}\frac{x_ix_j}{(1+ax_i)(1+ax_j)}.
\]
Each canceled product has nonnegative coefficients, so evaluation at $a>0$ bounds its degree-$d$ coefficient and proves~\eqref{eq:marked}.

The crucial summability constant is
\begin{equation}\label{eq:T}
 T:=\sum_{\substack{n\ge3\\n\text{ odd}}}e_2(x(n))
 \le\frac{\pi^2(4+3\sqrt2)}8<11.
\end{equation}
To verify it, repeat the absolute mixed-word argument at $m=1,d=2$, where the atomic eigenfunction is $1$. Now $Ag\le g/\sqrt2$, $W_11=2g$, and $W_1g\le\sqrt2g$. Summing the two zero-degree gaps bounds the absolute mixed sum by $2\sqrt2/(1-1/\sqrt2)^2$. Its orbit expression is exactly $(32/\pi^2)\sum_{n\ge1,\ n\text{ odd}}e_2(x(n))$, including the nearest term. This proves~\eqref{eq:T}.

At the saddle $a_*$, set $q=(1+\theta a_*)/(1+a_*)<1$. Equations~\eqref{eq:strictproduct}--\eqref{eq:T} show that the entire absolute remote mixed tail in the normalized tangent scale is at most
\begin{equation}\label{eq:tail}
 \frac{d^2}{d-1}a_*^2T C^dq^d=O(dC^dq^d).
\end{equation}

\section{The positive diagonal asymptotic}
The function $aB'(a)/B(a)=\sum_jab_j/(1+ab_j)$ increases continuously from $0$ to infinity, so $a_*$ exists uniquely. Let $X$ be the countable sum of independent Bernoulli variables with success probabilities $p_j=a_*b_j/(1+a_*b_j)$. Its mean is $1$, variance is $\sigma^2\in(0,\infty)$, and probability generating function is $B(a_*z)/B(a_*)$. The sum is finite almost surely, and all exponential moments exist because $B$ is entire. If $S_d$ is a sum of $d$ independent copies, then
\[
 g_d:=[a^d]B(a)^d=C^d\Pr(S_d=d).
\]
The elementary lattice local limit argument gives
\begin{equation}\label{eq:lclt}
 \Pr(S_d=d)\sim\frac1{\sqrt{2\pi d\sigma^2}}.
\end{equation}
For completeness, $X$ assigns positive mass to both $0$ and $1$, so its characteristic function has modulus strictly below $1$ away from $2\pi\Z$. Near zero, the centered characteristic function is $1-\sigma^2u^2/2+O(u^3)$. In Fourier inversion split at a small fixed neighborhood of zero. Outside it the integral decays exponentially in $d$. Inside it, the substitution $u=v/\sqrt d$, a Gaussian majorant, and dominated convergence give~\eqref{eq:lclt}.

The nearest mixed coefficient is $g_d-\sum_jb_j^d$, where $\sum_jb_j^d<2$ for $d\ge2$. Also $C>2$, using only the first two factors. The remote mixed contribution in~\eqref{eq:tail} is exponentially smaller than~\eqref{eq:lclt}, and $P_m=O((2/\pi)^d)$ by~\eqref{eq:purebound}. Substitution in~\eqref{eq:decomposition}, pairing positive and negative half-integers, gives~\eqref{eq:asymptotic}.

\section{An effective analytic cutoff}
The integer-mean mode theorem for countable Bernoulli sums states that the integer mean is a mode; see Gnedin's extension of Darroch's rule, Theorem~2 in \cite{gnedin}. It applies to $S_d$, whose mean is $d$ and variance is at most $d$. Chebyshev therefore gives
\[
 \Pr(S_d=d)\ge\frac{3/4}{4\sqrt d+1}\ge\frac1{8\sqrt d},
 \qquad g_d\ge\frac{C^d}{8\sqrt d}.
\]
Here the interval $|S_d-d|<2\sqrt d$ contains at most $4\sqrt d+1$ integers and has probability at least $3/4$.

The exact half-angle certificate in the package proves
\[
 \frac45<a_*<\frac56.
\]
It encloses the first twelve $b_j$ with rational bounds using $b_{j+1}=b_j/(1+\sqrt{1+b_j^2})$, and bounds the remaining mean tail geometrically. Specifically it proves the upper bound $981/1000$ for the mean at $a=4/5$ and the lower bound $1008/1000$ at $a=5/6$. Since $\theta<81/100$ and $q$ decreases with $a_*$,
\[
 q<\frac{206}{225}.
\]
The tail in~\eqref{eq:tail}, divided by $C^d$, is consequently below
\[
 \frac{275}{36}(d+2)\left(\frac{206}{225}\right)^d.
\]
For $d\ge140$ this is less than $1/(16\sqrt d)$. At $d=140$ the exact rational certificate checks
\[
 \left(\frac{275\cdot16}{36}\right)^2 142^2\cdot140
 \left(\frac{206}{225}\right)^{280}<1.
\]
The expression $(d+2)\sqrt d(206/225)^d$ decreases thereafter: its consecutive ratio is at most $(143/142)(281/280)(206/225)<1$.

Finally $\sum_jb_j^d<2<C^d/(16\sqrt d)$ for $d\ge140$, using $C>2$ and $2^d>32d$. The nearest mixed coefficient therefore strictly exceeds the absolute remote mixed tail. The positive $P_m$ proves $H_{m,m}>0$ for $m\ge70$.

At degree $4m$, the feedback formula is exactly
\begin{equation}\label{eq:pressure}
 E_{m,m}=\sum_{j=0}^mH_{m,j}[t^{4m}]\tan^{2m+2j}t
                  -\frac{R_m^2+R_{2m}}2.
\end{equation}
The preceding triangle supplies $H_{m,j}>0$ for $j<m$. The $j=0$ term is at least $R_m\binom{2m}{m}/3^m\ge2R_m/3$; the binomial ratio starts at $2/3$ for $m=2$ and its successive ratio is $(4m+2)/(3m+3)\ge1$. In contrast, $R_{2m}\le R_m^2$ and $R_m\le R_2=1/3$, so the two negative terms total at most $R_m/3$. Thus $H_{m,m}\ge0$ implies $E_{m,m}>0$, completing the analytic part of Theorem~\ref{thm:main} for $m\ge70$.

\section{Exact certification of the finite complement}\label{sec:finite}
The package uses a fixed-grid enclosure computation whose arithmetic consists solely of exact integers. It never decides a sign from an unbounded numerical approximation. Let $D=2^{2m-1}$, with Fourier indices $1-m,\ldots,m-1$. Define integer matrices
\[
 (B_j)_{k,\ell}=\binom{2m}{m+2k-\ell}
 [z^j](1-z)^{m+2k-\ell}(1+z)^{m-2k+\ell}.
\]
An entry is defined as zero when $m+2k-\ell$ lies outside $[0,2m]$. Then $V_a=D^{-1}\sum_j\ii^jB_ja^j$. The atomic vector is the Eulerian row of order $2m-1$, divided by $(2m-1)!$. Its normalization and eigenvector equation are checked exactly.

Write $h_a=\sum_n\ii^nv_na^n$. Below order $2m$, the eigenvalue has no nonconstant term, so
\[
 (DI-B_0)v_n=\sum_{j=1}^nB_jv_{n-j},\qquad\boldsymbol1^{\mathsf T}v_n=0.
\]
At $n=2m$, subtract $D(-1)^mR_mv_0$ from the right side. Reflection gives $(v_n)_{-k}=(-1)^n(v_n)_k$, permitting separate even and odd blocks. In the even block, the zero-index equation is replaced by normalization with row $(1,2,\ldots,2)$; in the odd block normalization is automatic.

Both block inverses are computed fraction-free, with every division checked exact and the final integer identity $CN=qI$ verified. The grid denominator is $S=2^{256}f^2$, where $f$ is the reduced denominator of $v_0$. This makes both $v_0$ and the feedback forcing exact on the grid. Every later inverse product is rounded downward to the grid, with the rounding inequality checked by integers.

Here is the rigorous error propagation. Use the full-vector $\ell^1$ norm on compressed coordinates: weight $1$ at index zero and weight $2$ at positive indices. Let $\alpha_n$ be the induced norm of the appropriate block inverse, excluding the normalization column in the even case. This is computed exactly from the integer inverse. Since
\[
 |B_j|\le\binom{2m}{j}B_0,\qquad
 \|B_j\|_1\le D\binom{2m}{j},
\]
an integer bound $U_n\ge S\|v_n-\widetilde v_n/S\|_1$ is propagated by
\[
 U_n=\left\lceil\alpha_nD\sum_{j=1}^n\binom{2m}{j}U_{n-j}\right\rceil
       +\begin{cases}2m-1,&n\text{ even},\\2m-2,&n\text{ odd},\end{cases}
 \qquad U_0=0.
\]
The last term bounds the total weighted rounding error. The half-period evaluation has norm one, so its computed value plus or minus $U_{2m}/S$ rigorously encloses $H_{m,m}$. All intermediate vectors and error bounds are saved in compressed replay traces. Independent full rational calculations and full-vector enclosures check selected cases.

All 68 moments $m=2,\ldots,69$ yield strictly positive rational lower endpoints. Their 4,828 response-order enclosures, baseline vectors, propagated errors, final intervals, and available full-rational comparisons were independently checked. This finite verification proves $H_{m,m}>0$ throughout the complement of the analytic range. Equation~\eqref{eq:pressure} then proves $E_{m,m}>0$, completing Theorem~\ref{thm:main}. The exact sources, every per-order vector, and the certificate checker are included; the finite computation is a distinct part of the proof.
% ed. (2026-10-01): the trace archives were not delivered; what the filed files show
\begin{ednote}
As filed, the per-order traces are missing: this package is the source archive of a split
delivery whose five trace archives (the 68 compressed files
\path{data/certificates/parity_mNNN.json.trace.gz}; see \path{DELIVERY_README.txt}) were not
delivered. The trace audit \path{code/check_trace_coverage.py} and the full replay
\path{code/reproduce_all.py} therefore cannot run on the filed files, and the baseline vectors,
the 4,828 propagated error bounds and the final vectors cannot be rechecked. What the filed
files show (checked on filing): the 68 interval records
\path{data/certificates/parity_mNNN.json} have strictly positive lower endpoints; exact
rational values of $H_{m,m}$ filed with the package lie inside them for the 35 orders
$2\le m\le35$ and $m=50$; for $m=68,69$ an independent full-vector enclosure overlaps them;
for $36\le m\le49$ and $51\le m\le67$ the filed intervals are the only finite evidence for
$H_{m,m}>0$. The pressure coefficient $E_{m,m}$ is proved again in \path{../Thue_Morse_Integer_Pressure_Full_Range/} (whose own finite intervals were not delivered either), and its positivity for
$2\le m\le128$ is certified by the filed intervals of \path{../Thue_Morse_Integer_Pressure_First_Negative_Data/}.
\end{ednote}
% ed. (2026-10-01): the missing traces can be reconstructed from the filed producer
\begin{ednote}
The missing traces can be regenerated, byte for byte, with the filed producer
\path{code/check_boundary_parity.cpp}: see the section ``Reconstructing the trace archives''
of \path{README.md}. On filing, the traces for $m=2,3,10,20,35,50,69$ were regenerated on a
copy: the interval records agree with the filed ones, the compressed traces reproduce the
hashes recorded in \path{data/trace_checks.json}, and the trace audit passes on them. The full
set (68 orders, about 45~MB compressed, a few minutes of computation) was not regenerated.
\end{ednote}

\section{Later negative coefficients and open directions}
The response theorem cannot uniformly extend to $r=m+1$. For $m=2$, set $x=a^2$ and write $\Lambda(a)=\lambda_2(a)$. The pressure cubic in \cite{integer}, under $\rho=\Lambda/(1+a^2)^2$ and $\cos2t=(1-x)/(1+x)$, gives
\[
 8\Lambda^3-2(7-x)(1+x)\Lambda^2
 +(7-3x)(1+x)^3\Lambda-(1+x)^6=0.
\]
Its derivative in $\Lambda$ at $(\Lambda,x)=(1,0)$ is $3$, so its local branch is unique. Exact substitution and a separate Fourier recursion agree on
\[
 \Lambda=1+\frac{x^2}{3}+8x^3+\frac{536}{27}x^4
             -\frac{488}{27}x^5-\frac{90472}{243}x^6+O(x^7).
\]
Since $\Lambda=1+a^4h_a(1/2)$, this proves $H_{2,3}=-488/27<0$. Pressure at the corresponding degree ten remains positive: $E_{2,3}=105448/2025$. Thus response and pressure sharpness are different questions.

A separate direct rational phase-eigenvalue recursion locates the first negative pressure degree \emph{above} $2m$ for the following small moments:
\[
\begin{array}{c|rrrrrrrrrrrrr}
 m&2&3&4&5&6&7&8&9&10&11&12&13&14\\\hline
 \text{degree}&12&18&24&32&38&44&52&58&64&72&78&84&92
\end{array}
\]
The lower even coefficients are already negative, so the qualification ``above $2m$'' matters. These are exact finite observations, not an all-order law. In particular, neither a limiting ratio nor a rounding formula for the first negative degree is asserted. Whether pressure alone stays positive throughout $1\le r<2m$ remains open here; the finite data are compatible with that extension. Positivity of each individual atomic-basis response coefficient throughout $r\le m$ is also a separate, unproved question.
% ed. (2026-10-01): pressure positivity below 6m and the first negative degree are settled by later packages
\begin{ednote}
Pressure positivity throughout $1\le r<2m$ is proved for every $m\ge2$ in
\path{../Thue_Morse_Integer_Pressure_Full_Range/}. The table above agrees with the certified first negative degrees
of \path{../Thue_Morse_Integer_Pressure_First_Negative_Data/} ($2\le m\le128$) and \path{../Thue_Morse_Integer_Pressure_Sign_Changes/}
($2\le m\le20$), and \path{../Thue_Morse_Integer_Pressure_First_Negative/} proves $N_m/m\to\gamma\in(6.662966,6.662968)$,
with a $\log\log m$ correction and a rounding rule valid away from the even lattice.
\end{ednote}

% ed. (2026-10-01): series map: the thirteen packages of the series and the sources cited here
\begin{ednote}
This note is the fourth of thirteen packages written in one series and filed together on 2026-10-01 beside the repository manuscript \path{../Thue_Morse_Integer_Pressure/} in the Fabius drafts tree (batch~72 of the repository's \path{docs/incoming/}; unreviewed). In logical order they are the positivity series \path{../Thue_Morse_Integer_Pressure_First_Positive/}, \path{../Thue_Morse_Integer_Pressure_Higher_Positive/}, \path{../Thue_Morse_Integer_Pressure_Positive_Triangle/}, \path{../Thue_Morse_Integer_Pressure_Feedback_Boundary/}, \path{../Thue_Morse_Integer_Pressure_Beyond_Boundary/}, \path{../Thue_Morse_Integer_Pressure_Linear_Region/} and \path{../Thue_Morse_Integer_Pressure_Full_Range/}, and the sign series, which imports the full-range theorem: \path{../Thue_Morse_Integer_Pressure_Sign_Changes/}, \path{../Thue_Morse_Integer_Pressure_Sign_Densities/}, \path{../Thue_Morse_Integer_Pressure_Negative_Bound/}, \path{../Thue_Morse_Integer_Pressure_First_Negative/} and \path{../Thue_Morse_Integer_Pressure_Cluster_Asymptotics/}, with the dataset \path{../Thue_Morse_Integer_Pressure_First_Negative_Data/} (no article). An earlier version of \path{../Thue_Morse_Integer_Pressure_Full_Range/}, \emph{The Full Pressure Positivity Range for Large Integer Orders} (orders $m\ge4096$), delivered in the same batch, is superseded by it and was not filed.
Here \cite{integer} is \path{../Thue_Morse_Integer_Pressure/} (the version cited differs from the filed manuscript only by dated editorial changes, none of them mathematical), and \cite{triangle} is
\path{../Thue_Morse_Integer_Pressure_Positive_Triangle/}.
\end{ednote}
\begin{thebibliography}{9}\small
\bibitem{integer}\emph{Integer Pressure and a Missing Taylor Coefficient}. ProveIt research manuscript, 28 September 2026. Inspected version: commit \nolinkurl{6a98f89bac85cc6b4ba5d5c3b63996037d9824e1}, blob \nolinkurl{1973fd17245864bb8a83e1b1d90416ac032ba7c1}.\newline
\href{https://github.com/VladimirReshetnikov/ProveIt/blob/6a98f89bac85cc6b4ba5d5c3b63996037d9824e1/Analysis/FabiusFunction/docs/semi-formalized-research-frontiers/drafts/thue-morse/Thue_Morse_Integer_Pressure/article.tex}{Versioned repository source}.
\bibitem{triangle}\emph{The Full Positive Triangle in Integer Thue Morse Pressure}. Separate research note prepared for Vladimir Reshetnikov with OpenAI, 1 October 2026. The tangent single-crossing theorem, full $r<m$ response and pressure positivity, and exact certificates.
\bibitem{gnedin} A. Gnedin, \emph{Cross Modality of the Extended Binomial Sums}, 2024, Theorem 2 (integer-mean mode rule for finite and countable Bernoulli sums).\newline\url{https://arxiv.org/html/2408.06477v1}
\end{thebibliography}
\end{document}
